\begin{opinion}	

La tesis de Josué presenta una metodología para resolver distintas variantes del VRP con poco tiempo de trabajo humano, aprovechando Algoritmos de Búsqueda Local. Lo más notable es que su enfoque elimina la necesidad de programar manualmente la evaluación de soluciones vecinas de casi cualquier variante, ya que este proceso se automatiza de manera eficiente en su propuesta. Este avance representa un paso importante en la optimización de recursos y tiempo en la investigación relacionada con el VRP.

Josué ha demostrado su capacidad para aplicar de manera creativa los conocimientos adquiridos durante su carrera. Ha demostrado una habilidad notable para identificar y resolver problemas emergentes en su investigación, destacándose en la adquisición autodidacta de conocimientos y habilidades necesarias para este proyecto. Su enfoque innovador ha sido fundamental en el éxito de su trabajo.

Estamos convencidos de que la tesis de Josué es un claro testimonio de su aptitud para desempeñarse como un destacado Científico de la Computación. Su investigación cumple con todos los requisitos de una tesis de licenciatura en este campo y pone de manifiesto su potencial para contribuir a la ciencia el futuro. Estamos seguros de que Josué continuará avanzando y enriqueciendo esta disciplina con sus futuras contribuciones.

\vspace{1cm}


\begin{flushright}
	\underline{\hspace{6.5cm}}\\
	MSc. Fernando Raúl Rodríguez Flores
	
	Facultad de Matemática y Computación
	
	Universidad de la Habana
	
	Noviembre, 2023
\end{flushright}

\end{opinion}